\section{Zerodivisors and total rings of fractions}
\label{section-total-quotient-ring}

\noindent
The local ring at a minimal prime has the following properties.

\begin{lemma}
\label{lemma-minimal-prime-reduced-ring}
Let $\mathfrak p$ be a minimal prime of a ring $R$.
Every element of the maximal ideal of $R_{\mathfrak p}$
is nilpotent. If $R$ is reduced then $R_{\mathfrak p}$
is a field.
\end{lemma}

\begin{proof}
If some element $x$ of ${\mathfrak p}R_{\mathfrak p}$
is not nilpotent, then $D(x) \not = \emptyset$, see
Lemma \ref{lemma-Zariski-topology}. This contradicts
the minimality of $\mathfrak p$. If $R$ is reduced,
then ${\mathfrak p}R_{\mathfrak p} = 0$ and
hence it is a field.
\end{proof}

\begin{lemma}
\label{lemma-reduced-ring-sub-product-fields}
Let $R$ be a reduced ring. Then
\begin{enumerate}
\item $R$ is a subring of a product of fields,
\item $R \to \prod_{\mathfrak p\text{ minimal}} R_{\mathfrak p}$
is an embedding into a product of fields,
\item $\bigcup_{\mathfrak p\text{ minimal}} \mathfrak p$ is the set
of zerodivisors of $R$.
\end{enumerate}
\end{lemma}

\begin{proof}
By Lemma \ref{lemma-minimal-prime-reduced-ring} each of the rings
$R_\mathfrak p$ is a field. In particular, the kernel of the ring
map $R \to R_\mathfrak p$ is $\mathfrak p$.
By Lemma \ref{lemma-Zariski-topology}
we have $\bigcap_{\mathfrak p} \mathfrak p = (0)$.
Hence (2) and (1) are true. If $x y = 0$ and $y \not = 0$, then
$y \not \in \mathfrak p$ for some minimal prime $\mathfrak p$.
Hence $x \in \mathfrak p$. Thus every zerodivisor of $R$ is contained
in $\bigcup_{\mathfrak p\text{ minimal}} \mathfrak p$.
Conversely, suppose that $x \in \mathfrak p$ for some minimal
prime $\mathfrak p$. Then $x$ maps to zero in $R_\mathfrak p$,
hence there exists $y \in R$, $y \not \in \mathfrak p$ such that
$xy = 0$. In other words, $x$ is a zerodivisor. This finishes the
proof of (3) and the lemma.
\end{proof}

\noindent
The total ring of fractions $Q(R)$ of a ring $R$ was introduced in
Example \ref{example-localize-at-prime}.

\begin{lemma}
\label{lemma-total-ring-fractions}
Let $R$ be a ring.
Let $S \subset R$ be a multiplicative subset consisting of nonzerodivisors.
Then $Q(R) \cong Q(S^{-1}R)$.
In particular $Q(R) \cong Q(Q(R))$.
\end{lemma}

\begin{proof}
If $x \in S^{-1}R$ is a nonzerodivisor, and
$x = r/f$ for some $r \in R$, $f \in S$, then
$r$ is a nonzerodivisor in $R$. Whence the lemma.
\end{proof}

\noindent
We can apply glueing results to prove something about
total rings of fractions $Q(R)$ which we introduced in
Example \ref{example-localize-at-prime}.

\begin{lemma}
\label{lemma-total-ring-fractions-no-embedded-points}
Let $R$ be a ring.
Assume that $R$ has finitely many minimal primes
$\mathfrak q_1, \ldots, \mathfrak q_t$, and that
$\mathfrak q_1 \cup \ldots \cup \mathfrak q_t$ is the set
of zerodivisors of $R$.
Then the total ring of fractions $Q(R)$ is equal to
$R_{\mathfrak q_1} \times \ldots \times R_{\mathfrak q_t}$.
\end{lemma}

\begin{proof}
There are natural maps $Q(R) \to R_{\mathfrak q_i}$ since
any nonzerodivisor is contained in $R \setminus \mathfrak q_i$.
Hence a natural map
$Q(R) \to R_{\mathfrak q_1} \times \ldots \times R_{\mathfrak q_t}$.
For any nonminimal prime $\mathfrak p \subset R$ we see that
$\mathfrak p \not \subset \mathfrak q_1 \cup \ldots \cup \mathfrak q_t$
by Lemma \ref{lemma-silly}. Hence
$\Spec(Q(R)) = \{\mathfrak q_1, \ldots, \mathfrak q_t\}$
(as subsets of $\Spec(R)$, see Lemma \ref{lemma-spec-localization}).
Therefore $\Spec(Q(R))$ is a finite discrete set and
it follows that $Q(R) = A_1 \times \ldots \times A_t$
with $\Spec(A_i) = \{\mathfrak{q}_i\}$, see
Lemma \ref{lemma-disjoint-implies-product}.
Moreover $A_i$ is a local ring, which is a localization
of $R$. Hence $A_i \cong R_{\mathfrak q_i}$.
\end{proof}

















